\par\medskip
\begin{center}\textls[100]{b) Akademische Tätigkeit in den letzten Lebensjahren.}\quad\textsc{Riemann.}\end{center}
\noindent\hspace*{\parindent}Über die letzten Lebensjahre lesen wir bei \textsc{Sartorius} (a. a. O. \textsuperscript{13)}, S. 67):
»Seine Vorlesungen über praktische Astronomie und über verschiedene Teile
der Mathematik, namentlich über die Methode der kleinsten Quadrate, schienen
ihm in vorgerückten Alter mehr Freude als in seinen früheren Lebensjahren
zu machen, und wenn er auch anfangs über diese ihm aufgebürdete Last beim
Beginn eines jeden neuen Semesters klagte, so war er doch bald mit grosser
Lebendigkeit und merkwürdiger geistiger Energie mitten in der Sache; auch
hat ein grösserer Kreis von Zuhörern, als der, den er früher um sich zu sehen\par
\clearpage\SourceLeaf{292}\ArithmeticHead{207}{R}{Xb. VORLESUNGSTÄTIGKEIT.}
\noindent gewohnt war, vielleicht anregend auf ihn gewirkt.« \textsc{R. Dedekind} gibt in dem
Aufsatze \textit{Gauss in seiner Vorlesung über die Methode der kleinsten Quadrate}\footnote{Festschrift, Göttingen 1901, S. 45—59, \textsc{Dedekinds} Werke II, S. 293—306.}
eine anschauliche Schilderung von \textsc{Gauss} und dem Milieu, in dem diese Vorlesung
stattfand, dieselbe übrigens, aus der Werke X, 1, S. 473—481 ein Auszug
nach einer von \textsc{August Ritter} herrührenden Ausarbeitung wiedergegeben
ist. \textsc{Dedekind} schreibt:
\par
»\mbox{.\,.\,.\,.\,.} Zu Ostern 1850 kam ich nach Göttingen; \mbox{.\,.\,.} auf meinen Wegen
nach oder von der Sternwarte \mbox{.\,.\,.} begegnete ich zuweilen \textsc{Gauss}, und erfreute
mich des Anblicks seiner stattlichen, Ehrfurcht gebietenden Erscheinung, und
sehr oft sah ich ihn in grösster Nähe auf seinem festen Platze im Literarischen
Museum, das er regelmässig besuchte, um Zeitungen zu lesen.
\par
Zu Anfang des folgenden Wintersemesters hielt ich mich für reif, seine
Vorlesung über die Methode der kleinsten Quadrate zu hören, und so betrat
ich, mit dem Testierbuch ausgerüstet und nicht ohne Herzklopfen, zum ersten
Male sein Wohnzimmer, wo ich ihn an seinem Schreibtisch sitzend fand.
Meine Meldung schien ihn wenig zu erfreuen, ich hatte auch wohl gehört,
dass er sich ungern entschloss Vorlesungen zu halten; nachdem er seinen
Namen in das Buch eingetragen hatte, sagte er nach kurzem Schweigen: Sie
wissen vielleicht, dass es immer sehr zweifelhaft ist, ob meine Vorlesungen
zustande kommen; wo wohnen Sie? bei dem Barbier \textsc{Vogel}? Nun, das trifft
sich ja glücklich, denn der ist auch mein Barbier, durch ihn werde ich Sie
benachrichtigen.
\par
Einige Tage darauf trat dann \textsc{Vogel}, eine stadtbekannte Persönlichkeit,
ganz erfüllt von der Wichtigkeit seiner Mission, bei mir ein, um zu bestellen,
dass sich noch mehr Zuhörer gemeldet hätten, und dass Herr Geh. Hofrat
\textsc{Gauss} die Vorlesung halten werde.
\par
Wir waren neun Studenten, von denen ich \textsc{A. Ritter} (später Professor
der Mechanik in Hannover) und \textsc{Moritz Cantor} (später Professor in Heidelberg)
nach und nach näher kennen lernte; \mbox{.\,.\,.} Das Auditorium war durch
ein Vorzimmer von \textsc{Gauss'} Arbeitszimmer getrennt und ziemlich klein. Wir
sassen an einem Tisch, dessen Längsseiten für je drei, aber nicht für vier
Personen bequemen Platz boten. Der Tür gegenüber, am obern Ende sass\par
\clearpage\SourceLeaf{293}\ArithmeticHead{208}{L}{L. SCHLESINGER, ÜBER GAUSS' FUNKTIONENTHEORETISCHE ARBEITEN.}
\noindent\textsc{Gauss} in mässiger Entfernung vom Tische, und wenn wir vollzählig waren,
so mussten zwei von uns, die zuletzt kamen, ganz in seine Nähe rücken und
ihr Heft auf den Schoss nehmen. \textsc{Gauss} trug ein leichtes schwarzes Käppchen,
einen ziemlich langen braunen Gehrock, graue Beinkleider; er sass
meist in bequemer Haltung, etwas gebeugt vor sich niedersehend, mit über
dem Leib gefalteten Händen. Er sprach ganz frei, sehr deutlich, einfach und
schlicht; wenn er aber einen neuen Gesichtspunkt hervorheben wollte, wobei
er ein besonders charakteristisches Wort gebrauchte, so erhob er wohl plötzlich
den Kopf, wandte sich zu seinem Nachbarn und blickte ihn während
der nachdrücklichen Rede ernst mit seinen schönen, durchdringenden blauen
Augen an. Das war unvergesslich. Seine Sprache war fast ganz dialektfrei,
nur bisweilen kamen Anklänge an unsere stadtbraunschweigische Mundart;
beim Zählen z. B., wobei er auch den Gebrauch der Finger nicht verschmähte,
sagte er nicht eins, zwei, drei, sondern eine, zweie, dreie usf., wie man es
noch jetzt bei uns auf dem Markte hören kann. Ging er von einer prinzipiellen
Erörterung zur Entwicklung mathematischer Formeln über, so erhob
er sich, und in stattlicher, ganz aufrechter Haltung schrieb er an einer neben
ihm stehenden Tafel mit der ihm eigenen schönen Handschrift, wobei es ihm
immer durch Sparsamkeit und zweckmässige Anordnung gelang, mit dem
ziemlich kleinen Raume auszukommen. Für die Zahlenbeispiele, auf deren
sorgfältige Durchführung er besonderen Wert legte, brachte er die erforderlichen
Data auf kleinen Zetteln mit\footnote{Es folgt nun eine Inhaltsübersicht über die wöchentlich dreistündige Vorlesung; vergl. \textsc{A. Galle} Werke XI, 2, Abh. 1, S. 6—7.}. \mbox{.\,.\,.} Ich kann nur sagen, dass wir
diesem ausgezeichneten Vortrage, in welchem auch einige Beispiele aus der
Theorie der bestimmten Integrale behandelt wurden, mit immer steigendem
Interesse gefolgt sind. Aber es schien uns auch, als ob \textsc{Gauss} selbst, der
vorher wenig Neigung gehabt hatte, die Vorlesung zu halten, im Laufe derselben
doch einige Freude an seiner Lehrtätigkeit empfand. So kam es am
13. März 1851 zum Schluss; \textsc{Gauss} erhob sich, wir alle mit ihm, und er entliess
uns mit den freundlichen Abschiedsworten: Es bleibt mir nur noch übrig,
Ihnen zu danken für die grosse Regelmässigkeit und Aufmerksamkeit, mit
der Sie meinem doch wohl recht trocken zu nennenden Vortrage gefolgt sind.
\mbox{.\,.\,.\,.}«\par
\clearpage\SourceLeaf{294}\ArithmeticHead{209}{R}{Xb. RIEMANNS PROMOTION UND HABILITATION.}
Aber auch über eine andere, nicht minder bedeutungsvolle Seite von \textsc{Gauss}
akademischer Tätigkeit aus diesen Jahren hat uns \textsc{Dedekind} unterrichtet, nämlich
über \textsc{Gauss'} Teilnahme an der Promotion und Habilitation \textsc{Bernhard
Riemanns}\footnote{Siehe \textsc{Riemanns} Werke, 2. Aufl. 1892, S. 539ff. insbesondere S. 545, 548, 549.}. \textsc{Riemanns} Dissertation \textit{Grundlagen für eine allgemeine Theorie der
Funktionen einer veränderlichen komplexen Grösse}, die im November 1851 der
Göttinger philosophischen Fakultät eingereicht worden war, fand, wie \textsc{Dedekind}
berichtet, eine sehr anerkennende Beurteilung durch \textsc{Gauss}, der \textsc{Riemann}
bei seinem Besuch mitteilte, dass er seit Jahren eine Schrift vorbereite, die
denselben Gegenstand behandele, sich aber freilich nicht darauf beschränke.
Unzweifelhaft hatte \textsc{Gauss} dabei die im vorhergehenden Abschnitt besprochene
Abhandlung über die Konvergenz der Reihen im Sinn, und für das lebhafte
Interesse, das \textsc{Gauss} der Erstlingsschrift seines genialen Zuhörers entgegen
brachte, zeugt ein Blättchen des Nachlasses (Fm, Varia analytica), wo unter
der Überschrift \textit{Taylors Lehrsatz nach Cauchys Behandlung} der \textsc{Cauchysche}
Integralsatz für den Fall eines Kreises abgeleitet wird, und auf dessen letzter
Seite der \textsc{Riemannsche} Integralsatz für das logarithmische Potential in der Form:
\[
-2\pi u_0=\int\left(u\frac{\partial\log r}{\partial p}-\log r\frac{\partial u}{\partial p}\right)ds
\]
steht. Die vollständige Übereinstimmung mit \textsc{Riemanns} Bezeichnungen (Inauguraldissertation,
Art. 10, \textsc{Riemanns} Werke, 2. Aufl., S. 20) bezeugt, dass
\textsc{Gauss} diese Formel sich aus \textsc{Riemanns} Inauguraldissertation notiert hat.
\par
Anfangs Dezember 1853 reichte \textsc{Riemann}, der inzwischen im\linebreak mathematisch-physikalischen
Seminar Assistent \textsc{Wilhelm Webers} geworden war, seine Habilitationsschrift
\textit{Über die Darstellbarkeit einer Funktion durch eine trigonometrische
Reihe} der Fakultät ein. Von den für die Probevorlesung vor der Fakultät
vorgeschlagenen drei Thematen wählte \textsc{Gauss}, gegen das übliche Herkommen,
nicht das erste, sondern das dritte \textit{Über die Hypothesen, welche der Geometrie
zu Grunde liegen}, »weil er begierig war zu hören, wie ein so schwieriger Gegenstand
von einem so jungen Manne behandelt werden würde«. »Nachdem ich«,
schreibt \textsc{Riemann} am 26. Juni 1854 an seinen Bruder, »etwa vierzehn Tage
nach Ostern mit einer andern Arbeit, die ich nicht gut vermeiden konnte,
fertig geworden war, ging ich nun eifrig an die Ausarbeitung meiner Probe-\par
\HistoricalFoot{X 2 Abh. 2.}{27}
\clearpage\SourceLeaf{295}\ArithmeticHead{210}{L}{L. SCHLESINGER, ÜBER GAUSS' FUNKTIONENTHEORETISCHE ARBEITEN.}
\noindent vorlesung und wurde um Pfingsten damit fertig. Ich erreichte es indess nur
mit vieler Mühe, dass ich mein Kolloquium gleich machen konnte. \mbox{.\,.\,.\,.} \textsc{Gauss'}
Gesundheitszustand ist nämlich in der letzten Zeit so schlimm geworden, dass
man noch in diesem Jahre seinen Tod fürchtet, und er sich zu schwach fühlte,
mich zu examinieren. \mbox{.\,.\,.} Ich hatte mich schon in das Unvermeidliche gefügt,
da entschloss er sich plötzlich auf mein wiederholtes Bitten, um die Sache
vom Halse los zu werden, am Freitag nach Pfingsten Mittag, das Kolloquium
auf den andern Tag um halb Elf anzusetzen und so war ich am Sonnabend
um Eins glücklich damit fertig.« Die Probevorlesung übertraf \textsc{Gauss'} Erwartungen
und setzte ihn in das grösste Erstaunen; auf dem Heimwege aus
der Fakultätssitzung sprach er sich zu \textsc{Wilhelm Weber} mit der höchsten Anerkennung
und mit einer bei ihm seltenen Erregung über die Tiefe der von
\textsc{Riemann} vorgetragenen Gedanken aus.
\par
Man darf es wohl als einen besonders harmonischen Ausklang dieses
reichen Lebens bezeichnen, dass es \textsc{Gauss} noch vergönnt war, die ersten
Schöpfungen des Mannes durch seinen Beifall zu krönen, der wie kein zweiter
berufen war, in der Funktionentheorie und in der Geometrie das, was \textsc{Gauss}
begonnen und erstrebt hatte, weiter zu fördern, durch seine tiefen und originalen
Gedanken zu befruchten und damit zum Ausgangspunkte aller weiteren
Entwicklung auf diesen Gebieten zu machen.
\par\bigskip\begin{center}\rule{.25\linewidth}{.3pt}\end{center}
